\section{Related Work}
\label{sec:related}

Principal component analysis has been applied to design evaluation in several domains.
In industrial design, Yanagisawa and Fukuda~\cite{yanagisawa2005study} used PCA to
identify independent aesthetic dimensions from multi-attribute product ratings.
In music, Aucouturier and Pachet~\cite{aucouturier2003music} applied PCA to timbral
features to identify orthogonal perceptual dimensions. In game studies, quantitative
approaches have been limited by corpus size: most studies work with fewer than 50 games
and rely on community ratings rather than structured axis scores.

The general factor problem we address --- where PCA on correlated evaluative ratings
yields a dominant quality PC --- is well-known in psychometrics as the \emph{g factor}
problem~\cite{spearman1904general} and in consumer research as the halo effect.
Within-group normalization to suppress the general factor has been used in
evolutionary biology~\cite{bookstein1991morphometric} and personality psychology
but has not, to our knowledge, been applied to game design analysis.
